\section{Materials and Methods}
\subsection{Study design, observations and environmental coverage}
We used the derived data and scripts released by Ortiz-Andrade and Rojas-Perez~\cite{ortiz2026}, together with the existing experiment outputs in this project. The study concerns Puerto Rico and retains complete eBird checklists using Stationary or Traveling protocols, durations of 5--240~min and travel distances no greater than 5~km. The response is one if the species was reported and zero otherwise. Zero denotes a non-report, not verified absence.

Habitat refers to the environment in which a species lives. We characterize selected habitat conditions through elevation and vegetation greenness and examine their associations with checklist records and modeled occupancy states. Habitat use is a broad description of these associations here; foraging, breeding behavior and habitat quality were not measured directly. At the observation level, a repeat report means that a checklist reports the species again at the same site in a later period. It depends on both presence and detection and is reported separately from modeled occupancy persistence.

The analysis covers 1 January 2013 to 31 December 2025. We retain the published period definitions. Pre-Mar\'ia ends on 19 September 2017; Inter-Hurricanes runs from 20 September 2017 to 17 September 2022; and Post-Fiona begins on 18 September 2022. These categories locate observations relative to the storms. They do not measure local storm exposure or supply an undisturbed counterfactual, so their contrasts are temporal associations rather than estimates of hurricane causation.

Elevation derives from SRTM~\cite{farr2007}; NDVI summarizes vegetation greenness~\cite{tucker1979}. The released Landsat procedure selects the nearest valid observation within $\pm16$ days and restricts imagery to the relevant period. We use the released covariates rather than repeating satellite extraction. The full analysis contains \nChecklists{} checklists, of which \nComplete{} have NDVI; \nMissingNDVI{} (\nMissingPct{}\%) do not. Environmental GLMs and the principal conventional benchmark use the complete subset. We describe missingness by period and reporting outcome without assuming it is random. This composition check is motivated by evidence that spatial and temporal sampling biases can interact in community-science trend estimation~\cite{backstrom2025}. The archived EVI2 comparison supplies a secondary vegetation-index check~\cite{jiang2008}.

Figure~\ref{fig:area} places checklist coverage and the validation blocks in the same landscape. Its 24,876 locations refer to the pre-extraction coordinate table; the exported environmental table contains 24,856 coordinate pairs because coordinate precision differs. Neither count is the number of occupancy sites or validation blocks. The sampling locations reflect opportunities for observation rather than a spatially uniform survey.
\begin{figure}[htbp]\centering
\incfig[0.94\linewidth]{fig01_study_area_and_spatial_blocks}
\caption{Puerto Rico, SRTM elevation and checklist coverage (upper panel), with one spatial fold allocation (lower panel). Entire $0.08^\circ$ blocks are assigned to folds; adjacent training and test blocks have no distance buffer.}\label{fig:area}
\end{figure}

\subsection{Verification and conditional reporting models}
\label{sec:verify}
File hashes, sample counts, filtering criteria and fitted coefficients were checked for data and model reproduction. Two duplicated date--coordinate records were retained to preserve the original analysis sample.
\ifzh
表~\ref{tab:screen}列出研究对象的来源筛选。参考研究的实施适合度得分为 5.1，高于其余候选的 3.4--4.1，主要优势是同时提供观测数据、环境协变量、分析脚本及校验信息。这使本研究能够从同一输入出发，连续检查报告格局、占域过程和预测验证，减少不同资料来源造成的比较差异。
\else
Table~\ref{tab:screen} summarizes source selection. The reference study scored 5.1 for implementation suitability, compared with 3.4--4.1 for the other candidates. Its combination of observations, environmental covariates, analysis scripts and checksums enabled a linked evaluation of reporting patterns, occupancy processes and predictive validation from a common input.
\fi
\inctab{screen}
\ifzh
数据核验通过 53/54 项检查，GLM 的 44/44 项检查均在规定容差内通过，11 个占域系数及标准误也得到复现（图~\ref{fig:verify}；表~\ref{tab:verify}）。唯一的数据检查差异涉及两条日期--坐标组合重复产生的多余记录，约占全部清单的 0.0022\%，且均未报告目标物种。该结果将数值差异的排查重点从文件版本和样本筛选转向模型内部的观测协变量对应关系。
\else
Data verification passed 53 of 54 checks; all 44 GLM checks met their tolerances, and all 11 occupancy coefficients and standard errors were reproduced (Figure~\ref{fig:verify}; Table~\ref{tab:verify}). The sole data-check discrepancy involved two excess records with duplicated date--coordinate combinations, approximately 0.0022\% of all checklists, both without a species report. This localized the subsequent investigation to observation-covariate alignment within the model rather than differences in input versions or filtering.
\fi
\begin{figure}[htbp]\centering\incfig{fig02_verification_and_reproduction}
\caption{\ifzh
数据核验与模型数值复现。颜色区分检查类别，系数比较使用原模型排列。
\else
Data verification and numerical model reproduction. Colors distinguish check categories; coefficient comparisons use the original model ordering.
\fi}\label{fig:verify}\end{figure}
\FloatBarrier
{\footnotesize\renewcommand{\arraystretch}{0.96}\inctab{verify}}
\FloatBarrier


Five binomial GLMs progress from an intercept through period, observation effort, environment and period--elevation interaction. All use the same NDVI-complete subset. The published selected model is
\begin{equation}
\operatorname{logit}(q_i)=\alpha+\beta_{P(i)}+\beta_E E_i+\delta_{P(i)}E_i+\beta_V V_i+\beta_D D_i+\beta_L L_i,
\label{eq:glm}
\end{equation}
where $q_i$ is reporting probability, $E_i$ elevation in 100-m units, $V_i$ NDVI, $D_i$ duration in minutes and $L_i$ distance in kilometres. Pre-Mar\'ia is the reference period. We report existing AIC, model weights and pseudo-$R^2$ as summaries of the candidate fits~\cite{burnham2002,mcfadden1974}, separately from held-out performance.

For the elevation curves, NDVI, duration and distance are fixed at their complete-sample means. Predictions at 100 and 900~m summarize low and high elevations, not biological thresholds. Following the released analysis, we also summarize the curve on a uniform 0--1300~m grid at 10-m intervals:
\begin{equation}
C_t=\frac{\sum_h h\widehat q_t(h)}{\sum_h\widehat q_t(h)}.
\end{equation}
This probability-weighted elevation centroid describes a fitted curve, not population abundance, geographic range centre or tracked movement. Annual reporting intervals reproduce the archived binomial intervals and do not account for observer or geographic dependence.

\subsection{Occupancy persistence and observation-model correction}
\label{sec:ordering}
The dynamic model uses $0.01^\circ$ rounded-coordinate cells as sites and the three periods as primary occasions. Sites require a checklist in every period. Within each site--period, the first three records after ordering by year and original row order are retained, with fewer visits represented by missing cells. The resulting \nOccSites{}-by-nine detection history contains \nOccasions{} observed occasions and \nOccMissing{} missing positions. Site covariates are computed from all eligible records before visit selection and standardized across retained sites. NDVI is therefore a cross-time site summary, not a within-site vegetation-recovery trajectory.

The selected \texttt{unmarked} model~\cite{fiske2011} relates initial occupancy ($\psi$) and local extinction ($\varepsilon$) to elevation and NDVI, colonization ($\gamma$) to NDVI, and detection ($p$) to period. Period does not separately enter the transition components. Thus the model does not estimate a storm-specific change in extinction or colonization.

Let $z_{g,t}$ denote the latent occupancy state of grid $g$ in primary period $t$. Occupancy persistence probability is defined as $\Pr(z_{g,t+1}=1\mid z_{g,t}=1)=1-\varepsilon_g$, the probability that a previously occupied grid remains occupied in the next period~\cite{mackenzie2003}. It does not represent residence of the same individual, individual survival or uninterrupted presence throughout a period. Colonization probability $\gamma$ conditions on prior nonoccupancy, whereas an observed first or repeat report does not directly identify colonization. Persistence throughout this paper refers to this grid-level state transition.

The observation covariate must align with the detection matrix in site-major order. All nine observations from one site precede those of the next. The released script instead repeats each occasion label across sites. The archived labelled-frame check and controlled refit change only this alignment, keeping the response, site covariates, formula and optimization settings fixed. We compare estimates and uncertainty, using $\widehat\beta\pm1.96\,\mathrm{SE}$ for approximate 95\% Wald intervals. This is a sensitivity test within the selected model, not a new selection across all corrected candidate models. Multi-year primary periods also place substantial demands on within-period closure; estimated transitions are not annual demographic rates.

\subsection{Reporting prediction in held-out geographic blocks}
\label{sec:protocols}
The prediction target remains a checklist report. The eleven conventional configurations are GLM m3, GLM m4, a spline GLM, m4 with terrain, random forest, LightGBM, XGBoost, a multilayer perceptron (MLP), and terrain-augmented variants of the last three. Tree methods provide flexible tabular predictors~\cite{breiman2001,ke2017,chen2016}. Their features include environment, effort, period, seasonal terms and coordinates; the published GLMs retain their original covariates. Terrain configurations add 20 SRTM descriptors of elevation, slope, aspect and neighbourhood variation at four scales. Consequently, comparisons concern implemented model--information combinations, not isolated architecture effects.

The saved settings use 400 trees and a minimum leaf size of five for random forest; the boosting implementations use 600 estimators and learning rate 0.05, with 63 leaves for LightGBM and maximum depth six for XGBoost. The MLP has three hidden layers of width 128. Estimated transformations, including MLP scaling and spline knots, are fitted using each training fold. These settings describe the existing runs; we do not retune against the reported tests.

The primary protocol assigns entire $0.08^\circ$ geographic blocks, approximately 8--9~km per side, to ten size-balanced folds, repeated with three block origins. The eleven configurations share all 30 test folds. Each fold can contain disconnected blocks; no buffer separates adjoining training and test blocks. The protocol tests the specified within-island geographic holdout, not arbitrary-distance extrapolation. Secondary random validation uses three stratified five-fold repetitions. Training fractions and test composition follow the respective partition rules. The archived temporal experiment repeats fits on the same Post-Fiona test sample with three seeds; it evaluates transfer of relationships learned in earlier periods to later reports. Separate geographic and temporal holdouts follow the principle that validation should represent the intended transfer task~\cite{spacetime2025}.

We emphasize AUC, average precision (AP), Brier score and recall among the highest-ranked 10\% of test checklists~\cite{saito2015,brier1950}. They address overall discrimination, positive-report ranking, probability error and a fixed ranking budget, respectively. Comparisons use unweighted fold means and between-fold standard deviations, absolute differences, relative improvements and top-decile report coverage. Training sets overlap across repeated partitions, so folds are not treated as independent replicates for significance classification. The pooled curves use a separate ten-fold prediction export. ROC and precision--recall characterize ranking, while reliability diagrams use 12 equal-count bins and expected calibration error uses 15. Scale sensitivity evaluates $0.04^\circ$, $0.08^\circ$, $0.16^\circ$ and $0.32^\circ$ blocks; the EVI2 substitution uses six folds.

\subsection{\ifzh 环境分层、过程概率与监测预算实验\else Environmental stratification, process probabilities and monitoring budgets\fi}
\label{sec:envmethods}
\ifzh
补充分析将海拔分为低于 300~m、300--不足 600~m 和不低于 600~m 三层；NDVI 采用 68,976 条完整清单的三分位切点 0.4172 和 0.7006，形成九个环境组合。这些分层用于探索环境异质性，不作为物种的生境阈值。占域分析将相同 NDVI 切点应用于地点跨时期均值，并报告各组合的地点数；少于十个地点的组合单独标记。
\else
Supplementary analyses stratified elevation as below 300~m, 300--below 600~m, and at least 600~m. NDVI tertile cutpoints of 0.4172 and 0.7006, computed from 68,976 complete checklists, produced nine environmental combinations. These exploratory strata describe heterogeneity rather than species habitat thresholds. Occupancy summaries apply the same NDVI cutpoints to cross-period site means, report site counts, and flag combinations with fewer than ten sites.
\fi
\ifzh
NDVI 缺失比较先汇总各海拔层的报告率与缺失率，再拟合缺失指示变量与报告之间的二项关联模型，依次加入时期和海拔样条，以及调查时长的对数、距离的 $\log(1+x)$ 变换、季度和经纬度二次项。标准误按 214 个 $0.08^\circ$ 空间区块聚类。模型使用信赖域优化获得稳定起点，并检查收敛、梯度及协方差。报告回升使用复现的 m4 公式重新拟合，将每个环境组合内全部记录的协变量分布固定，分别替换时期后求平均预测。Fiona 后相对 María 前的差值区间通过聚类协方差和 delta 法计算；原始分层报告率另进行 1,000 次空间区块自助抽样。
\else
NDVI-missing comparisons first summarize report and missingness rates within elevation bands. Binomial reporting models then relate reports to a missingness indicator, sequentially adding period and an elevation spline, followed by log duration, $\log(1+x)$ distance, quarter, and quadratic geographic-coordinate terms. Standard errors cluster 214 geographic blocks of width $0.08^\circ$. Trust-region optimization supplies stable starting values, with convergence, gradient and covariance checks. Reporting rebound is standardized using a refit of the reproduced m4 formula. Each environmental cell retains its pooled covariate distribution while period is replaced in turn. Post-Fiona minus Pre-Mar\'ia intervals use clustered covariance and the delta method; raw stratified rates additionally use 1,000 spatial block bootstrap samples.
\fi
\ifzh
校正占域模型在 DGX 上重新拟合，系数与已有校正结果的最大差异小于 $10^{-4}$。对每个地点预测初始占域、占域持续概率 $1-\varepsilon$ 和定殖 $\gamma$，再按环境组合平均。2,000 次联合正态参数模拟保留完整协方差，给出概率与对比的 95\% 区间。固定海拔 100~m、NDVI 从 0.30 增至 0.76 的对比用于分离低海拔内的绿度响应；这些取值近似对应低海拔低、高绿度地点的中位数。另从九列检测历史中筛选相邻两个时期各有三次访问的地点--时期对，统计前期已报告后再次报告的比例。此比例描述等访问次数下的重复报告，与潜在占域转移分别呈现。
\else
The corrected occupancy model was refitted on the DGX, reproducing existing corrected coefficients within $10^{-4}$. Site predictions of initial occupancy, occupancy persistence probability $1-\varepsilon$, and colonization $\gamma$ were averaged by environmental cell. Two thousand joint normal coefficient draws preserve the full covariance and yield 95\% probability and contrast intervals. A contrast at 100~m from NDVI 0.30 to 0.76 isolates the fitted greenness response; these values approximate the low- and high-greenness medians among low-elevation sites. Adjacent period pairs with three visits in each period were also extracted from the nine-column detection history to summarize repeated reports conditional on a previous report. This equal-visit reporting proportion is presented separately from latent occupancy transitions.
\fi
\ifzh
新增随机森林采用原 400 棵树、最小叶节点样本数五和种子零，在一次十折空间划分中重新训练。逐折核验事件标识符和地理区块的训练--测试交集均为空，并重新计算 m4 的首折预测验证原预测文件的行对应关系。分层预测比较包括 m4、随机森林和两种地形提升树。监测预算在每个测试折固定为清单数的 10\% 向下取整。全局方案直接选该折最高分记录，分层方案按各海拔层样本量分配约 10\%，采用最大余数分配保持每折总名额完全相同。两方案均选择 6,890 条清单。报告 AP、Brier、各层入选比例及正报告覆盖率；覆盖率差的区间采用 1,000 次空间区块自助抽样，条件于本次划分的拟合预测。该实验评价已有清单的排序与环境覆盖，不把清单数直接换算为实地访问成本。
\else
A new random forest used the original 400 trees, minimum leaf size five and seed zero in one ten-fold spatial partition. Every fold was checked for disjoint training/test event identifiers and geographic blocks; recomputing the first m4 fold validated the row alignment of the saved predictions. Stratified comparisons include m4, random forest and both terrain-augmented boosting models. Each test fold receives a budget equal to the floor of 10\% of its checklist count. Global ranking selects the highest scores within that fold; stratified ranking allocates approximately 10\% within each elevation band, using largest remainders to preserve exactly the same fold budget. Both select 6,890 checklists. AP, Brier, selection fractions and positive-report coverage are reported by stratum. Coverage-contrast intervals use 1,000 spatial block bootstrap samples conditional on this partition's fitted predictions. The experiment evaluates record ranking and environmental coverage; checklist counts are not converted into field-visit costs.
\fi


Figure~\ref{fig:arch} summarizes how the environmental and survey data support complementary analyses of conditional reporting, occupancy persistence and reporting prediction, and how these analyses inform follow-up monitoring across elevation and greenness conditions.
\begin{figure}[htbp]\centering\incfig[\linewidth]{主图/1}
\caption{\ifzh
环境分析与监测建议概览。（a）eBird 记录、SRTM 海拔、NDVI 和调查协变量用于条件报告、动态占域与随机森林报告预测三类互补分析。（b）分析结果为固定山地复访、低地绿度对照及环境分层内候选记录排序提供依据，以指导后续调查。
\else
Overview of the environmental analyses and monitoring implications. (a) eBird records, SRTM elevation, NDVI and survey covariates inform complementary analyses of conditional reporting, dynamic occupancy and random forest reporting prediction. (b) The findings motivate fixed upland revisits, lowland greenness contrasts and within-stratum candidate ranking to guide follow-up surveys.
\fi}\label{fig:arch}
\end{figure}
\FloatBarrier

\FloatBarrier

\subsection{Neural prediction structure and training-source diagnostics}
\label{sec:audit}
The archived HERON workflow factorizes a reporting score into two bounded branch scores, $q_i=s_i d_i$, and combines it with conventional models. We reconstruct the archived fold assignments and match sampling-event identifiers between the complete and all-record tables. The original workflow partitions these tables independently but treats equal fold numbers as equal held-out records. We assess overlap with the neural training source before interpreting its performance; internal stacking uses an analogous independent allocation. The diagnostic calculates overlap from the original allocation rules and observation identifiers.

Separately, the grouped loss uses $1-\prod_i(1-s_i d_i)$, which depends only on the products. It therefore does not by itself identify either branch as occupancy or detection. For a shared site state, an occupancy model instead has $\Pr(\hbox{at least one report})=\psi_g[1-\prod_j(1-p_{gj})]$. Input restrictions and monotone effort curves do not establish equivalence to this shared-state likelihood. Accordingly, neural environmental and effort branches are interpreted as fitted scores, while occupancy persistence is evaluated through the dynamic occupancy transitions. Neural-component ablations use 20 folds, stack ablations use six, and the direct configuration comparison uses ten; contrasts are analyzed within their respective fold sets.

